\renewcommand{\bibname}{Quellen und Bandzuordnung}
\begin{thebibliography}{9}
\addcontentsline{toc}{chapter}{Quellen und Bandzuordnung}
\footnotesize
\setlength{\itemsep}{4pt}
\setlength{\parskip}{0pt}
\interlinepenalty=10000

\bibitem{project}
Martin Kunze: \emph{Die Grundlagen der Mathematik}.
Band~6: Injektivität, Bildmengen, Einschränkungen,
Zielmengenwechsel, Fallabbildungen und Kompositionen.
Band~3 und Band~5 liefern die verwendeten Mengen- und
Funktionsbegriffe. Format nach den Leitfäden zu
Band~3, Band~4, Band~5, Band~28 und Band~37.

\bibitem{lafferriere}
Beatriz Lafferriere, Gerardo Lafferriere und Mau Nam Nguyen:
\emph{Introduction to Mathematical Analysis I}.
3.~Auflage, Portland State University, 2022.
Abschnitt~1.2, Definition~1.2.2(ii): Injektivität;
Satz~1.2.5(ii),(iv): Komposition und Rückschluß auf den
inneren Faktor; Übungen~1.2.1(a) und~1.2.4(a),(b):
Urbild einer Bildmenge, Schnitte und Differenzen.
\url{https://content.library.pdx.edu/files/PDXScholar/Math/Introduction-to-Mathematical-Analysis-I-3rd-Edition/}.

\bibitem{campesato}
Jean-Baptiste Campesato:
\emph{MAT246H1-S, Homework Questions Week 10},
University of Toronto, 29.~März bis 2.~April 2021.
Aufgabe~6: Charakterisierung der Injektivität durch
die Erhaltung binärer Durchschnitte unter Bildbildung.
\url{https://www.math.utoronto.ca/campesat/ens/21W/q10.pdf}.

\bibitem{ucl}
University College London:
\emph{MATH0005 Algebra 1}, Abschnitt~2.10,
\emph{Conditions for invertibility}.
Satz~2.10.1(1): Charakterisierung der Injektivität durch
eine Linksinverse bei nichtleeren Mengen.
Hier wird ausschließlich dieser Injektivitätsteil verwendet.
\url{https://www.homepages.ucl.ac.uk/~ucahmto/0005_2021/Ch2.S10.html}.

\end{thebibliography}

\noindent\textbf{Einordnung der Zusammenstellung.}
Die Quellen dienen der Kontrolle der Begriffe und Kernaussagen.
Die konkreten Zusammenfassungen, die systematische Behandlung
leerer Mengen und die ergänzenden mengenbasierten Konstruktionen
sind eigenständig formuliert. Nicht jede Aussage steht in den
Quellen in genau dieser Gestalt. Der Leitfaden beansprucht weder
historische Erstentdeckungen noch eine vollständige Erfassung
aller bekannten Resultate. Er enthält ausschließlich Definitionen,
Sätze und kurze Hinweise zu ihrem Geltungsbereich.